\chapter{Introduction}

For as long as I can remember, I have wanted to formalize the concepts that govern how people relate to one another.
Not to drain them of feeling, but to give them precision.
Relationships and behavior are domains where language tends to go soft at exactly the wrong moment --- where the most important things are said in ways that leave too much room for misreading.
I have been misunderstood enough times to believe that the gap between what someone means and what someone hears is not merely a failure of attention, but a structural problem: the words we use to describe relational phenomena are often underspecified, carrying different meanings to different people with no clear way to resolve the difference.

Mathematics offers a remedy.
A well-formed definition does not shift depending on who is reading it.
A proposition either holds or it does not.
This does not mean that human relationships are simple --- it means that when we describe them carefully, we can at least agree on what we are describing.
The goal of this book is to build that kind of careful description: a formal language for patterns of relationship and behavior that recur across everyday life.

A secondary motivation is classification.
Patterns that feel vague or ineffable often turn out, once written down precisely, to be instances of something general.
Giving a pattern a name and a definition is not a reduction of its significance --- it is a tool for recognizing it, reasoning about it, and communicating it to someone else without losing anything in translation.

The ultimate intent of this work is practical: to offer readers more clarity in their own relationships than they had before.
Toward that end, the book is organized around propositions --- formally stated, but chosen because they capture something true about how relationships actually work.
Many of these propositions will be familiar as intuition.
The claim that actions carry more weight than words, for instance, is something most people already believe.
What formal treatment adds is not the insight itself, but a precise account of what the statement actually means, under what conditions it holds, and what follows from it.
That kind of precision, it is hoped, is useful not just as an intellectual exercise but as something a person can take into their own life.
